\documentclass[10pt,a4paper]{article}

\usepackage[utf8]{inputenc}
\usepackage[T1]{fontenc}

\usepackage{amsmath,mathtools,amsfonts,amssymb}
\usepackage{graphicx}
\usepackage[spanish]{babel}
\usepackage{lastpage,tabto,xcolor}
\usepackage[a4paper, top=0.8in, bottom=1in, left=0.7in, right=0.7in]{geometry}
\usepackage{fancyhdr}
\usepackage{comment}
\usepackage{multicol}
\usepackage{titling} % Para ajustar los espacios del título
\usepackage{tasks} % Para listas en columnas
\usepackage{float}
\usepackage[mode=image]{standalone}
\usepackage{tikz}
\usepackage{pgfplots}
\pgfplotsset{width=7cm,compat=1.15}
\usepackage{caption}

% Fecha de versión automática según fecha de compilación
\newcommand{\verdate}{\the\year.\ifnum\month<10 0\fi\the\month.\ifnum\day<10 0\fi\the\day}

% Ajusta la separación entre el encabezado y el contenido
\setlength{\headsep}{10pt}

% Ajuste del título:
\setlength{\droptitle}{-0.5in} % Mueve el título ligeramente hacia arriba
\pretitle{\begin{center}\vspace{-1em}\normalfont\Large}
	\posttitle{\par\end{center}\vspace{-4em}}
\preauthor{\begin{center}\vspace{-0.5em}}
	\postauthor{\par\end{center}}
\predate{}\postdate{}

\pagestyle{fancy}
\lhead{\footnotesize Análisis de Señales y Sistemas  - Año \the\year{} Ver. \verdate}
\rhead{\footnotesize Trabajo Práctico N$^\text{o}$6: Señales discretas y muestreo.}
\rfoot{\footnotesize Página \thepage\ de \pageref{LastPage}}
\renewcommand{\headrulewidth}{0.4pt}
\renewcommand{\footrulewidth}{0.4pt}

% Título optimizado con espacios reducidos
\title{%
	{\normalsize Departamento de Ingeniería en Tecnologías Electrónicas, UTN FRM}\\[0.1cm]
	{\large Análisis de Señales y Sistemas}\\[0.15cm]
	{\normalsize Trabajo Práctico N$^{\text o}$6: Señales de tiempo discreto y teorema del muestreo}%
}
\author{}
\date{}


\begin{document}
	\maketitle
	\thispagestyle{fancy}

	% Entorno de dos columnas
	\begin{multicols}{2}

		\label{sec:enun}

		{\small\noindent En todo el trabajo, $\Delta t$ es el intervalo de
		muestreo, $\omega_s = 2\pi/\Delta t$ la frecuencia de muestreo y
		$\lambda = \omega\,\Delta t$ la frecuencia digital, medida en radianes
		por muestra.}

		\vspace{0.4em}

		% ***********************************************
		% 		PARTE I -- SEÑALES DE TIEMPO DISCRETO
		% ***********************************************
		\textbf{Del mundo continuo al digital}
		\begin{enumerate}

		% ***********************************************
		% 				Ejercicio 1
		% ***********************************************
		\item Una tensión senoidal $x(t) = 1{,}2\cos(2\pi\cdot 100\,t)$ V se
		muestrea cada $\Delta t = 1{,}25$ ms y se cuantiza después con cuatro
		bits por muestra en signo-magnitud, con un paso de grilla de $0{,}2$ V.
		\begin{enumerate}
			\item Obtener $x[n]$ y simplificar el argumento del coseno hasta
			dejarlo en la forma $\cos(\lambda n)$.
			\item Indicar cuántas muestras entran en un ciclo de $x(t)$ y cuál
			es el período fundamental $N$ de la secuencia.
			\item Verificar que la grilla de niveles cubre todo el rango de la
			señal, e indicar cuántos niveles distintos quedan definidos.
			\item Completar una tabla con la muestra exacta, la magnitud
			entera, la palabra de cuatro bits y el valor cuantizado, para
			$n = 0, 1, \dots, 7$.
			\item Indicar el error máximo de redondeo.

		\end{enumerate}

		% ***********************************************
		% 		La secuencia: anatomía y operaciones
		% ***********************************************
		\textbf{La secuencia: anatomía y operaciones}

		% ***********************************************
		% 				Ejercicio 2
		% ***********************************************
		\item Sea la secuencia
		\[
			x[n] = \{\, -1 \quad 2 \quad \underset{\uparrow}{3} \quad 1 \quad -2 \,\}.
		\]
		\begin{enumerate}
			\item Graficar $x[n]$ e indicar su soporte.
			\item Graficar $x[n-3]$, $x[-n]$ y $x[2-n]$, indicando en cada caso
			el soporte resultante.
			\item Indicar cuáles de las cuatro secuencias anteriores son
			causales.
		\end{enumerate}

		% ***********************************************
		% 				Ejercicio 3
		% ***********************************************
		\item Sea la secuencia
		\[
			x[n] = \{\, \underset{\uparrow}{1} \quad -1 \quad 2 \quad 0 \quad 1 \quad 3 \,\}.
		\]
		\begin{enumerate}
			\item Determinar y graficar la decimación $y[n] = x[2n]$.
			\item Determinar y graficar la interpolación $w[n] = x[n/3]$.
			\item Indicar si se puede recuperar $x[n]$ a partir de $y[n]$, y si
			se puede a partir de $w[n]$. Justificar en los dos casos.
			\item Explicar por qué el escalado temporal $x(at)$ del tiempo continuo no
			se traslada al discreto como una sola operación.
		\end{enumerate}

		% ***********************************************
		% 				Ejercicio 4
		% ***********************************************
		\item Dada la secuencia
		\[
			x[n] = \{\, 6 \quad 4 \quad \underset{\uparrow}{-2} \quad -2 \,\},
		\]
		\begin{enumerate}
			\item Determinar sus componentes par e impar, volcando el cálculo
			en una tabla índice por índice.
			\item Graficar las tres secuencias.
			\item Verificar que $x_i[0] = 0$ y que la suma de las componentes
			reconstruye $x[n]$.
		\end{enumerate}

		% ***********************************************
		% 	La exponencial compleja y la frecuencia digital
		% ***********************************************
		\textbf{La exponencial compleja discreta y la frecuencia digital}

		% ***********************************************
		% 				Ejercicio 5
		% ***********************************************
		\item Para cada una de las frecuencias digitales siguientes, hallar la
		frecuencia equivalente en el intervalo $(-\pi,\pi]$:
		\begin{center}
		$\lambda_1 = \dfrac{\pi}{3}$, \quad
		$\lambda_2 = \dfrac{5\pi}{3}$, \quad
		$\lambda_3 = \dfrac{7\pi}{3}$, \quad
		$\lambda_4 = \pi$, \quad
		$\lambda_5 = \dfrac{11\pi}{6}$, \quad
		$\lambda_6 = 3\pi$.
		\end{center}
		\begin{enumerate}
			\item Ordenarlas de menor a mayor velocidad de oscilación.
			\item Indicar qué pares producen exactamente la misma secuencia
			$e^{j\lambda n}$ y qué pares comparten solo la parte real.
			\item Justificar, con la cuenta, por qué sumarle $2\pi$ a $\lambda$
			deja la secuencia intacta.
		\end{enumerate}

		% ***********************************************
		% 				Ejercicio 6
		% ***********************************************
		\item Un vector giratorio de frecuencia $\omega = 2\pi\cdot 50$ rad/s se
		muestrea cada $\Delta t$ segundos.
		\begin{enumerate}
			\item Calcular $\lambda$ para $\Delta t = 1$ ms, $6$ ms, $15$ ms y
			$20$ ms.
			\item Para cada caso, determinar la frecuencia digital equivalente
			en $(-\pi,\pi]$ y describir el giro aparente: qué ángulo avanza el
			vector entre foto y foto y, cuando ese giro existe, en qué sentido
			ocurre.
			\item Determinar para qué valores de $\Delta t$ el giro puede
			seguirse sin ambigüedad, y analizar qué pasa en el caso límite.
		\end{enumerate}

		% ***********************************************
		% 		Periodicidad y señales armónicas
		% ***********************************************
		\textbf{Periodicidad y señales armónicas}

		% ***********************************************
		% 				Ejercicio 7
		% ***********************************************
		\item Determinar si las siguientes secuencias son periódicas. Cuando lo
		sean, hallar el período fundamental $N$ e indicar, para las que son una
		sola exponencial o sinusoide, cuántas vueltas $k$ da el vector giratorio
		en ese período.
		\begin{enumerate}
			\item $x[n] = \cos(\pi n/5)$
			\item $x[n] = e^{j3\pi n/4}$
			\item $x[n] = \cos(15 n)$
			\item $x[n] = \cos(5\pi n)$
			\item $x[n] = \cos(\pi n/3) + \sin(\pi n/4)$
		\end{enumerate}

		% ***********************************************
		% 				Ejercicio 8
		% ***********************************************
		\item La señal $x(t) = 3\cos(20t)$ se muestrea con un intervalo
		$\Delta t$ elegido de modo que la secuencia $x[n]$ resulte periódica con
		período fundamental $N = 10$.
		\begin{enumerate}
			\item Determinar todos los valores posibles de $\lambda$ y de
			$\Delta t$ dentro de una vuelta.
			\item Para cada solución, indicar cuántos ciclos de $x(t)$
			transcurren mientras la secuencia recorre un período.
			\item Señalar cuáles de esas soluciones muestrean por encima del
			doble de la frecuencia de la señal, y relacionar esa condición con
			la posición de $\lambda$ respecto de $\pi$.
		\end{enumerate}

		% ***********************************************
		% 				Ejercicio 9
		% ***********************************************
		\item Se toma $N = 6$ y el conjunto de armónicas
		$\phi_k[n] = e^{j(2\pi/6)kn}$.
		\begin{enumerate}
			\item Escribir las armónicas distintas y verificar con la cuenta que
			$\phi_6[n] = \phi_0[n]$.
			\item Las armónicas $\phi_4$ y $\phi_5$ tienen frecuencia entre
			$\pi$ y $2\pi$. Hallar para cada una la frecuencia equivalente en
			$(-\pi,0)$, verificar sobre un período que la secuencia es la misma e
			indicar qué relación guarda cada una con otra armónica del conjunto.
			Señalar cuál de las seis oscila más rápido.
			\item Comprobar que $\lambda = 2\pi/5$ y $\lambda = 4\pi/5$ tienen
			el mismo período $N = 5$ pero distinta velocidad de oscilación.
		\end{enumerate}

		% ***********************************************
		% 		La exponencial general: la secuencia z^n
		% ***********************************************
		\textbf{La exponencial general: la secuencia $z^n$}

		% ***********************************************
		% 				Ejercicio 10
		% ***********************************************
		\item Sea $z = r\,e^{j\lambda}$ y la secuencia causal
		$x[n] = z^n\,u[n]$.
		\begin{enumerate}
			\item Para $z = 0{,}9\,e^{j\pi/6}$, escribir
			$\operatorname{Re}\{x[n]\}$, identificar la envolvente y el período
			de la oscilación, y graficar las primeras seis muestras.
			\item Repetir para $z = 1{,}1\,e^{j\pi/6}$ y comparar las dos
			gráficas.
			\item Para $z = -0{,}8$, mostrar que $z^n = (0{,}8)^n(-1)^n$ e
			indicar qué valor de $\lambda$ corresponde a ese punto del plano.
			\item Clasificar las secuencias de los incisos (a) y (b) como de
			energía, de potencia o de ninguna de las dos.
			\item Indicar qué puntos del plano complejo hacen que
			$\operatorname{Re}\{x[n]\}$ sea una sinusoide discreta de amplitud
			constante.
		\end{enumerate}

		% ***********************************************
		% 	Escalón, impulso y representación por impulsos
		% ***********************************************
		\textbf{Escalón, impulso y representación por impulsos}

		% ***********************************************
		% 				Ejercicio 11
		% ***********************************************
		\item Sean las secuencias
		\begin{align*}
			x_1[n] &= u[n+2] - u[n-3], \\
			x_2[n] &= u[n+1] + u[n-1] - u[n-3] - u[n-5], \\
			x_3[n] &= n\,\big(u[n] - u[n-5]\big), \\
			x_4[n] &= 2\,\delta[n+1] - \delta[n] + 3\,\delta[n-2].
		\end{align*}
		\begin{enumerate}
			\item Graficar las cuatro secuencias.
			\item Escribir $x_3[n]$ como suma de impulsos desplazados.
		\end{enumerate}

		% ***********************************************
		% 		PARTE II -- TEOREMA DEL MUESTREO
		% ***********************************************
		\textbf{El problema del muestreo}

		% ***********************************************
		% 				Ejercicio 12
		% ***********************************************
		\item Una rueda gira a $20$ vueltas por segundo. Sobre el borde lleva
		pintada una marca, que es lo único que permite reconocer en qué
		posición angular está la rueda en cada cuadro. Se la filma a $24$
		cuadros por segundo.
		\begin{enumerate}
			\item Calcular el ángulo que avanza la marca entre cuadro y cuadro.
			\item Reducir ese ángulo al intervalo $(-\pi,\pi]$ y determinar a
			qué velocidad, y en qué sentido, parece girar la rueda en la
			pantalla.
			\item Determinar a qué velocidades de giro la rueda parece quieta.
			\item Determinar la máxima velocidad de giro que esa filmación
			representa sin ambigüedad.
		\end{enumerate}

		% ***********************************************
		% 		El espectro de una señal muestreada
		% ***********************************************
		\textbf{El espectro de una señal muestreada}

		% ***********************************************
		% 				Ejercicio 13
		% ***********************************************
		\item Una señal $x(t)$ tiene espectro $X(j\omega)$ triangular, de altura
		$1$ en el origen y nulo fuera de $|\omega| < 3000\pi$ rad/s.
		\begin{enumerate}
			\item Con $\Delta t = 0{,}2$ ms, calcular $\omega_s$ y graficar
			$X_s(j\omega)$ mostrando al menos tres copias, con sus posiciones
			y su factor de escala. Demarcar la banda base y verificar sobre la
			gráfica que $X_s(j\omega)$ es periódico de período $\omega_s$.
			\item Repetir con $\Delta t = 0{,}4$ ms e indicar en qué banda de
			frecuencias se produce el solapamiento.
			\item Repetir las dos gráficas con el eje medido en frecuencia
			digital $\lambda = \omega\,\Delta t$, indicando en cada caso cuánto
			ocupa una copia y dónde queda la banda base. Escribir la condición
			de Nyquist en términos de $\lambda$.

		\end{enumerate}

		% ***********************************************
		% 				Nyquist y aliasing
		% ***********************************************
		\textbf{Nyquist y aliasing}

		% ***********************************************
		% 				Ejercicio 14
		% ***********************************************
		\item Para cada una de las siguientes señales, determinar la frecuencia
		máxima $\omega_M$ y la velocidad de Nyquist, y proponer una frecuencia de
		muestreo con margen sobre esa velocidad, indicando el intervalo $\Delta t$
		correspondiente.
		\begin{enumerate}
			\item $x(t) = 10\cos(5t) - 5\sin(7\pi t) + 3\cos(30t)$
			\item $y(t) = 18\sin(10t) + 15\sin(5\pi t) + 5\cos(60t)$
			\item $z(t) = 15\cos(6\pi t) + 7\sin(10\pi t) + 2\cos(15\pi t)$

		\end{enumerate}

		% ***********************************************
		% 				Ejercicio 15
		% ***********************************************
		\item Para la señal $z(t)$ del ejercicio anterior se adopta
		$\omega_s = 80$ rad/s.
		\begin{enumerate}
			\item Verificar si se cumple la condición de Nyquist.
			\item Indicar qué componentes quedan bien muestreadas.
			\item Calcular la frecuencia alias de la componente mal muestreada
			e indicar dónde queda respecto de las componentes bien muestreadas.
			\item Indicar si esa componente puede recuperarse con un
			procesamiento posterior. Justificar.
			\item Determinar la mínima frecuencia de muestreo que evita el
			problema.
		\end{enumerate}

		% ***********************************************
		% 				Ejercicio 16
		% ***********************************************
		\item El espectro de amplitud de una señal $x(t)$ es el de la
		Figura~\ref{fig:tp6_espectro_trapecio}, nulo fuera de $|\omega| < 500$
		rad/s. Se la muestrea con una frecuencia de muestreo $f_s = 150$ Hz.

		\begin{center}
			\includestandalone[width=0.98\linewidth]{figures/fig_tp6_espectro_trapecio/fig_tp6_espectro_trapecio}
			\captionof{figure}{Espectro de amplitud del ejercicio 16.}
			\label{fig:tp6_espectro_trapecio}
		\end{center}

		\begin{enumerate}
			\item Calcular $\omega_s$ y la frecuencia de Nyquist, demarcar la
			banda base y verificar si se cumple la condición de no
			solapamiento.
			\item Determinar qué porción del espectro se solapa y sobre qué
			banda de la banda base se deposita.
			\item Graficar el espectro muestreado sobre la banda base,
			señalando la zona contaminada.
			\item Analizar cómo se ve la componente de $\omega = 460$ rad/s
			después del muestreo.
			\item Determinar la mínima $f_s$ que evita el solapamiento.
		\end{enumerate}

		% ***********************************************
		% 				Del modelo a la práctica
		% ***********************************************
		\textbf{Del modelo a la práctica}

		% ***********************************************
		% 				Ejercicio 17
		% ***********************************************
		\item Una señal útil ocupa la banda $|\omega| < \omega_{\max}$ y se le
		suma un ruido concentrado en $\omega_1 < |\omega| < \omega_2$, según la
		Figura~\ref{fig:tp6_senal_ruido}. Tomar
		$\omega_{\max} = 1000$ rad/s, $\omega_1 = 1500$ rad/s y
		$\omega_2 = 2000$ rad/s.

		\begin{center}
			\includestandalone[width=0.98\linewidth]{figures/fig_tp6_senal_ruido/fig_tp6_senal_ruido}
			\captionof{figure}{Espectro de la señal útil (A) y del ruido (B),
			ejercicio 17.}
			\label{fig:tp6_senal_ruido}
		\end{center}

		\begin{enumerate}
			\item Con $\omega_s = 2500$ rad/s, demarcar la banda base, ubicar
			las réplicas del ruido y determinar si contaminan la banda útil.
			\item Proponer dos maneras distintas de muestrear sin afectar la
			banda útil, y para cada una indicar qué condición debe cumplir
			$\omega_s$.
			\item Explicar por qué el filtro antialiasing debe ser analógico y
			actuar antes del muestreo.
			\item Explicar por qué en la práctica la frecuencia de corte se
			ubica algo por debajo de la de Nyquist.
		\end{enumerate}

		% ***********************************************
		% 	El eje de frecuencias medido en lambda
		% ***********************************************
		\textbf{El eje de frecuencias medido en $\lambda$}

		% ***********************************************
		% 				Ejercicio 18
		% ***********************************************
		\item Una señal de audio abarca de $20$ Hz a $20$ kHz y se muestrea
		con $f_s = 44{,}1$ kHz. Una señal de vibración mecánica abarca de
		$0{,}5$ Hz a $500$ Hz.
		\begin{enumerate}
			\item Ubicar los extremos de la banda de audio sobre el eje
			$\lambda$ y demarcar la banda base.
			\item Determinar la frecuencia de muestreo que hace que la señal
			de vibración ocupe exactamente el mismo tramo del eje $\lambda$,
			y el intervalo $\Delta t$ correspondiente.
			\item Representar sobre el eje $\lambda$ la banda que ocupa cada
			una de las dos señales muestreadas y comparar las dos
			representaciones.
			\item Calcular, en Hz, la frecuencia continua que corresponde a
			$\lambda = \pi/3$ y a $\lambda = \pi$ en cada uno de los dos
			casos, y explicar por qué el mismo $\lambda$ representa
			frecuencias tan distintas.
		\end{enumerate}

		\end{enumerate}

	\end{multicols}

\end{document}
